% !TeX root = ../main.tex
\chapter{PHÂN TÍCH VÀ THIẾT KẾ HỆ THỐNG}
\label{ch:thietke}

\section{Phân tích yêu cầu}
\label{sec:yeucau}

\subsection{Yêu cầu chức năng}

\begin{table}[htbp]
\centering
\caption{Danh sách yêu cầu chức năng}
\label{tab:req-cn}
\small
\begin{tabular}{c p{9.5cm}}
\toprule
\textbf{Mã} & \textbf{Yêu cầu chức năng}\\
\midrule
FN1 & Hệ thống cung cấp dịch vụ định danh tập trung dựa trên Keycloak.\\
FN2 & Người dùng đăng nhập một lần (SSO) vào ứng dụng demo qua giao thức OIDC.\\
FN3 & Quản trị viên tạo và quản lý realm, client, role, user dưới dạng mã nguồn.\\
FN4 & Hệ thống hỗ trợ xác thực đa yếu tố (TOTP).\\
FN5 & Cấu hình định danh được áp dụng tự động qua quy trình CI/CD.\\
FN6 & Ứng dụng demo hiển thị thông tin người dùng nhận từ token.\\
\bottomrule
\end{tabular}
\end{table}

\subsection{Yêu cầu phi chức năng}

\begin{table}[htbp]
\centering
\caption{Danh sách yêu cầu phi chức năng}
\label{tab:req-pcn}
\small
\begin{tabular}{c p{9.5cm}}
\toprule
\textbf{Mã} & \textbf{Yêu cầu phi chức năng}\\
\midrule
NF1 & Bảo mật: sử dụng HTTPS/TLS, luồng Authorization Code + PKCE, không lưu mật khẩu dạng thô.\\
NF2 & Tự động hóa: có thể tái lập toàn bộ hệ thống từ mã nguồn.\\
NF3 & Khả năng bảo trì: cấu hình tách khỏi image, quản lý bằng Git.\\
NF4 & Hiệu năng: đáp ứng đăng nhập với độ trễ chấp nhận được ở quy mô thử nghiệm.\\
NF5 & Khả năng mở rộng: kiến trúc sẵn sàng chạy nhiều bản sao và tự phục hồi.\\
\bottomrule
\end{tabular}
\end{table}

\subsection{Tác nhân và ca sử dụng}
\label{subsec:usecase}

Bảng~\ref{tab:usecase} liệt kê các tác nhân và ca sử dụng chính của hệ thống.

\begin{table}[htbp]
\centering
\caption{Tác nhân và ca sử dụng chính}
\label{tab:usecase}
\small
\begin{tabular}{p{2.6cm} p{4.2cm} p{6.4cm}}
\toprule
\textbf{Tác nhân} & \textbf{Ca sử dụng} & \textbf{Mô tả}\\
\midrule
Người dùng & Đăng nhập SSO & Xác thực qua Keycloak và truy cập ứng dụng demo\\
Người dùng & Bật MFA & Thiết lập và sử dụng mã TOTP\\
Quản trị viên & Quản lý realm/client & Tạo, sửa cấu hình định danh dưới dạng mã nguồn\\
Quản trị viên & Triển khai & Áp dụng cấu hình qua pipeline CI/CD\\
Hệ thống CI & Kiểm thử và build & Chạy kiểm thử, build và đẩy image\\
ArgoCD & Đồng bộ & Đồng bộ trạng thái mong muốn từ Git lên cụm\\
\bottomrule
\end{tabular}
\end{table}

\section{Kiến trúc tổng thể}
\label{sec:kientruc}

Hệ thống được thiết kế theo kiến trúc ba tầng như minh họa trên
Hình~\ref{fig:kientruc}. Tầng dưới cùng cung cấp dịch vụ định danh;
tầng giữa chịu trách nhiệm quản lý cấu hình định danh như mã nguồn;
tầng trên cùng đảm nhiệm tự động hóa hạ tầng và vận hành.

\begin{figure}[htbp]
\centering
\begin{tikzpicture}[
  font=\small,
  layer/.style={draw, rounded corners, minimum width=13cm, align=center,
                inner sep=7pt}
]
\node[layer, fill=blue!8] (l3)
  {\textbf{TẦNG 3 -- TỰ ĐỘNG HÓA HẠ TẦNG}\\[2pt]
   CI/CD (GitHub Actions) $\cdot$ IaC (Terraform/Ansible) $\cdot$
   GitOps (ArgoCD)\\
   Kubernetes (k3s) $\cdot$ Giám sát (Prometheus/Grafana)};
\node[layer, fill=green!8, below=5mm of l3] (l2)
  {\textbf{TẦNG 2 -- CẤU HÌNH ĐỊNH DANH NHƯ MÃ NGUỒN}\\[2pt]
   Realm $\cdot$ Client $\cdot$ Role $\cdot$ User $\cdot$ Chính sách MFA\\
   (keycloak-config-cli / Keycloak Terraform Provider)};
\node[layer, fill=orange!12, below=5mm of l2] (l1)
  {\textbf{TẦNG 1 -- DỊCH VỤ ĐỊNH DANH}\\[2pt]
   Keycloak (OIDC/OAuth2, SSO, MFA) $+$ PostgreSQL\\
   $+$ Ứng dụng demo (Relying Party)};
\end{tikzpicture}
\caption{Kiến trúc tổng thể ba tầng của hệ thống}
\label{fig:kientruc}
\end{figure}

Việc tách thành ba tầng giúp phân định rõ trách nhiệm: tầng dịch vụ tập trung
vào logic định danh, tầng cấu hình đảm bảo tính nhất quán và có thể truy vết
của cấu hình, còn tầng tự động hóa chịu trách nhiệm triển khai và vận hành.

\subsection{Mô hình kiến trúc C4}
\label{subsec:c4}

Để mô tả kiến trúc một cách có hệ thống, chuyên đề sử dụng mô hình C4 ở hai
mức: sơ đồ ngữ cảnh (context) và sơ đồ container. Ở mức ngữ cảnh, hệ thống gồm
ba nhóm tác nhân: người dùng cuối, quản trị viên và các hệ thống tự động
(CI/CD, GitOps). Người dùng tương tác với ứng dụng demo; ứng dụng demo tin cậy
Keycloak để xác thực; quản trị viên và hệ thống tự động tương tác với Keycloak
và hạ tầng qua API.

Ở mức container (Hình~\ref{fig:c4}), hệ thống gồm các thành phần: ứng dụng demo
đóng vai trò Relying Party, Keycloak (IdP), PostgreSQL lưu trữ, cơ chế cấu hình
như mã nguồn, cụm Kubernetes (k3s), ArgoCD, pipeline CI và hệ thống giám sát.

\begin{figure}[htbp]
\centering
\begin{tikzpicture}[font=\small,
  box/.style={draw, rounded corners, minimum height=0.9cm, align=center, inner sep=4pt},
  arr/.style={-{Latex[length=2mm]}, thick}]
\node[box, fill=blue!8] (user) {Người dùng};
\node[box, fill=orange!12, right=1.5cm of user] (app) {Ứng dụng demo\\\textit{(Relying Party)}};
\node[box, fill=green!10, right=1.5cm of app] (kc) {Keycloak\\\textit{(IdP)}};
\node[box, fill=gray!10, below=0.8cm of kc] (db) {PostgreSQL};
\node[box, fill=blue!8, below=1.0cm of app] (cicd) {CI/CD \& GitOps};
\node[box, fill=blue!8, below=1.0cm of user] (admin) {Quản trị viên};
\draw[arr] (user) -- node[above,font=\scriptsize]{HTTPS} (app);
\draw[arr] (app) -- node[above,font=\scriptsize]{OIDC} (kc);
\draw[arr] (kc) -- (db);
\draw[arr] (admin) -- (kc);
\draw[arr] (cicd) -- (kc);
\end{tikzpicture}
\caption{Sơ đồ container của hệ thống (mức C4)}
\label{fig:c4}
\end{figure}

\section{Thiết kế luồng xác thực OIDC}
\label{sec:luongoidc}

Luồng xác thực trong hệ thống tuân theo quy trình Authorization Code kèm PKCE:

\begin{enumerate}
  \item Người dùng truy cập ứng dụng demo và chọn ``Đăng nhập''.
  \item Ứng dụng chuyển hướng người dùng tới Keycloak kèm
  \texttt{code\_challenge}.
  \item Keycloak xác thực người dùng (mật khẩu và, nếu được bật, mã TOTP).
  \item Keycloak chuyển hướng trở lại ứng dụng kèm \texttt{authorization code}.
  \item Ứng dụng gửi \texttt{authorization code} và \texttt{code\_verifier} tới
  endpoint \texttt{/token} để đổi lấy ID Token và Access Token.
  \item Ứng dụng xác thực token và thiết lập phiên cho người dùng.
\end{enumerate}

% TODO: Vẽ biểu đồ tuần tự (sequence diagram) minh họa 6 bước trên.

Bảng~\ref{tab:oidcflow} mô tả chi tiết từng bước của luồng, bao gồm tác nhân
và dữ liệu trao đổi.

\begin{table}[htbp]
\centering
\caption{Các bước của luồng Authorization Code kèm PKCE}
\label{tab:oidcflow}
\small
\begin{tabular}{c p{2.8cm} p{9.4cm}}
\toprule
\textbf{Bước} & \textbf{Tác nhân} & \textbf{Hành động và dữ liệu}\\
\midrule
1 & Người dùng $\rightarrow$ Ứng dụng & Chọn đăng nhập.\\
2 & Ứng dụng $\rightarrow$ Keycloak & Chuyển hướng tới \texttt{/authorize} kèm \texttt{code\_challenge}, \texttt{state}.\\
3 & Keycloak $\rightarrow$ Người dùng & Hiển thị trang đăng nhập và xác thực (mật khẩu, TOTP).\\
4 & Keycloak $\rightarrow$ Ứng dụng & Chuyển hướng về \texttt{redirect\_uri} kèm \texttt{authorization code}.\\
5 & Ứng dụng $\rightarrow$ Keycloak & Gửi code và \texttt{code\_verifier} tới \texttt{/token}.\\
6 & Keycloak $\rightarrow$ Ứng dụng & Trả về ID Token, access token, refresh token.\\
7 & Ứng dụng $\rightarrow$ Người dùng & Xác thực token và thiết lập phiên.\\
\bottomrule
\end{tabular}
\end{table}

\section{Thiết kế quản lý cấu hình định danh như mã nguồn}
\label{sec:cfgascode}

Toàn bộ cấu hình định danh được mô tả trong kho Git, gồm:
\begin{itemize}
  \item Định nghĩa realm và các thuộc tính bảo mật (thời hạn token,
  chính sách mật khẩu).
  \item Định nghĩa client cho ứng dụng demo (loại confidential/public,
  redirect URI, phạm vi scope).
  \item Định nghĩa role và ánh xạ vai trò.
  \item Cấu hình MFA (TOTP).
\end{itemize}

Khi cấu hình thay đổi trong Git, một pipeline sẽ áp dụng thay đổi đó vào
Keycloak, đảm bảo môi trường luôn nhất quán và không phụ thuộc thao tác thủ
công trên giao diện quản trị.

\section{Thiết kế quy trình triển khai tự động}
\label{sec:quytrinh}

Quy trình triển khai được thiết kế như sau:

\begin{enumerate}
  \item Lập trình viên đẩy thay đổi lên kho Git.
  \item Hệ thống CI (GitHub Actions) chạy kiểm thử, xây dựng image ứng dụng và
  đẩy lên registry.
  \item Các manifest Kubernetes được cập nhật trong kho Git (hoặc kho riêng
  cho cấu hình).
  \item ArgoCD phát hiện thay đổi và đồng bộ trạng thái mong muốn lên cụm k3s.
  \item Hệ thống giám sát (Prometheus/Grafana) theo dõi trạng thái dịch vụ.
\end{enumerate}

Cơ chế này cho phép triển khai theo kiểu cuốn chiếu (rolling update) và tự
động khôi phục khi có sự cố, hạn chế gián đoạn dịch vụ.

\section{Thiết kế cơ sở dữ liệu và trạng thái phiên}
\label{sec:csdl}

Keycloak lưu trữ dữ liệu realm, người dùng, client và phiên đăng nhập trong cơ
sở dữ liệu quan hệ (PostgreSQL). Bên cạnh đó, phiên và dữ liệu realm được đệm
trong bộ nhớ Infinispan để tăng tốc truy cập. Bảng~\ref{tab:data} tóm tắt các
nhóm dữ liệu chính.

\begin{table}[htbp]
\centering
\caption{Các nhóm dữ liệu chính trong hệ thống}
\label{tab:data}
\small
\begin{tabular}{p{3.4cm} p{3.2cm} p{7cm}}
\toprule
\textbf{Nhóm dữ liệu} & \textbf{Lưu trữ} & \textbf{Mô tả}\\
\midrule
Cấu hình realm/client & PostgreSQL & Định nghĩa realm, client, role\\
Tài khoản người dùng & PostgreSQL & Thông tin người dùng và mật khẩu đã băm\\
Phiên đăng nhập & Infinispan + cookie & Trạng thái SSO của người dùng\\
Token & Không lưu (stateless) & Được sinh và xác thực theo chữ ký\\
\bottomrule
\end{tabular}
\end{table}

\section{Thiết kế tính sẵn sàng cao}
\label{sec:ha}

Để đảm bảo dịch vụ định danh luôn sẵn sàng, kiến trúc được thiết kế với nhiều
bản sao Keycloak đặt sau bộ cân bằng tải, cơ sở dữ liệu có sao chép và phiên
được đệm phân tán. Các điểm cuối kiểm tra sức khỏe (liveness/readiness) giúp
Kubernetes tự phát hiện và thay thế Pod lỗi, trong khi cơ chế cuốn chiếu đảm
bảo cập nhật không gián đoạn. Bộ tự động mở rộng (HPA) giúp tăng số bản sao khi
tải tăng.

\section{Thiết kế bảo mật}
\label{sec:thietkebaomat}

Trên cơ sở các mối đe dọa đã phân tích ở Chương~1, Bảng~\ref{tab:secdesign}
ánh xạ từng nhóm mối đe dọa với biện pháp thiết kế tương ứng.

\begin{table}[htbp]
\centering
\caption{Ánh xạ mối đe dọa và biện pháp thiết kế}
\label{tab:secdesign}
\small
\begin{tabular}{p{3.4cm} p{9.4cm}}
\toprule
\textbf{Mối đe dọa} & \textbf{Biện pháp thiết kế}\\
\midrule
Nghe lén / MITM & Bắt buộc TLS cho mọi kết nối\\
Đánh cắp phiên/token & Thời hạn token ngắn, PKCE, cookie an toàn\\
Dò mật khẩu & Giới hạn tần suất, chính sách mật khẩu, MFA\\
Truy cập trái phép & RBAC, đặc quyền tối thiểu, NetworkPolicy\\
Lộ bí mật & Quản lý bí mật bằng Secret/Vault, không nhúng vào mã nguồn\\
Sửa đổi cấu hình & Quản lý cấu hình qua Git, kiểm soát phiên bản\\
\bottomrule
\end{tabular}
\end{table}

\section{Thiết kế cấu trúc kho mã nguồn}
\label{sec:repo}

Kho mã nguồn được tổ chức thành hai phần: mã nguồn ứng dụng (kèm Dockerfile) và
mã nguồn vận hành (manifest Kubernetes, cấu hình realm, pipeline CI, khai báo
GitOps). Việc tách nhánh theo môi trường (dev, staging, prod) và sử dụng pull
request giúp kiểm soát thay đổi trước khi áp dụng. Bảng~\ref{tab:repo} mô tả
cấu trúc thư mục dự kiến.

\begin{table}[htbp]
\centering
\caption{Cấu trúc thư mục kho mã nguồn dự án}
\label{tab:repo}
\small
\begin{tabular}{p{3.6cm} p{9.2cm}}
\toprule
\textbf{Thư mục} & \textbf{Nội dung}\\
\midrule
\texttt{app/} & Mã nguồn ứng dụng demo và Dockerfile\\
\texttt{deploy/} & Tệp docker-compose và cấu hình realm\\
\texttt{k8s/} & Manifest Kubernetes (Deployment, Service, Ingress, HPA)\\
\texttt{gitops/} & Khai báo ArgoCD và các overlay theo môi trường\\
\texttt{.github/} & Pipeline CI/CD\\
\texttt{docs/} & Tài liệu, sơ đồ thiết kế\\
\bottomrule
\end{tabular}
\end{table}

\section{Thiết kế quản lý bí mật}
\label{sec:secrets}

Thông tin nhạy cảm như mật khẩu cơ sở dữ liệu, khóa ký token và mật khẩu quản
trị tuyệt đối không được lưu trong mã nguồn. Ở môi trường thử nghiệm, chúng
được đưa vào qua biến môi trường hoặc tệp không được phiên bản hóa; ở môi
trường triển khai, chúng được quản lý bằng Kubernetes Secret, Sealed Secrets
hoặc HashiCorp Vault. Nguyên tắc chung là tách bí mật khỏi mã và chỉ cấp cho
thành phần cần dùng.

\section{Thiết kế thành phần hệ thống}
\label{sec:thanhphan}

Hệ thống gồm các thành phần chính: (1) ứng dụng demo -- giao diện người dùng và
xử lý luồng OIDC; (2) Keycloak -- dịch vụ định danh trung tâm; (3) PostgreSQL
-- lưu trữ dữ liệu định danh; (4) công cụ cấu hình như mã nguồn; (5) pipeline
CI; (6) ArgoCD; và (7) hệ thống giám sát. Mỗi thành phần có trách nhiệm riêng
và giao tiếp qua các giao diện xác định, giúp hệ thống dễ thay thế và mở rộng.

\section{Kết chương}

Chương này đã phân tích yêu cầu chức năng, phi chức năng và thiết kế kiến trúc
ba tầng của hệ thống, cùng các luồng xác thực và quy trình triển khai tự động.
Chương tiếp theo trình bày quá trình triển khai thực tế dựa trên thiết kế này.
